\chapter{Infraestrutura de TI}

\section{Escopo do edital e visão de auditor}

O edital de Auditor/TI não cobra infraestrutura como uma coleção de produtos. O conteúdo percorre redes TCP/IP, gerenciamento, LAN/WAN, segurança de redes, criptografia, sistemas operacionais, diretórios, virtualização, servidores de aplicação, alta disponibilidade, data center, armazenamento, backup, contêineres e DevOps. A preparação deve conectar três planos: \textbf{funcionamento técnico}, \textbf{risco operacional} e \textbf{evidência de auditoria}.

\begin{teoria}[modelo mental]
Para qualquer componente de infraestrutura, responda:
\begin{enumerate}
\item qual serviço ele presta e em que camada atua;
\item de quais componentes depende;
\item como falha e qual é seu domínio de falha;
\item quais controles reduzem probabilidade ou impacto;
\item quais métricas e logs demonstram que o controle opera;
\item como recuperar o serviço e os dados após falha.
\end{enumerate}
Esse raciocínio transforma conhecimento de administração de sistemas em conhecimento de Auditor/TI.
\end{teoria}

\section{Redes de computadores e arquitetura TCP/IP}

\subsection{Camadas, encapsulamento e unidades de dados}

O modelo TCP/IP separa responsabilidades para permitir interoperabilidade. Uma aplicação produz dados; o transporte adiciona informação de multiplexação e, no caso do TCP, confiabilidade; a camada de Internet adiciona endereçamento e roteamento IP; a camada de acesso encapsula o pacote em quadro apropriado ao enlace.

\begin{center}
\begin{tikzpicture}[node distance=4mm]
\tikzset{bl/.style={draw=azulPetroleo,rounded corners,minimum width=9cm,minimum height=.9cm,align=center}}
\node[bl,fill=azulClaro] (a) {Aplicação: HTTP(S), DNS, SMTP, SSH, SNMP, NTP};
\node[bl,fill=white,below=of a] (t) {Transporte: TCP / UDP};
\node[bl,fill=azulClaro,below=of t] (i) {Internet: IPv4 / IPv6 / ICMP};
\node[bl,fill=white,below=of i] (e) {Acesso: Ethernet, 802.11, VLAN};
\draw[-Latex,ambar,thick] (a.east) -- ++(1,0) |- (e.east) node[midway,right,align=left]{encapsulamento};
\draw[-Latex,azulMedio,thick] (e.west) -- ++(-1,0) |- (a.west) node[midway,left,align=right]{desencapsulamento};
\end{tikzpicture}
\end{center}

Um switch de camada 2 aprende endereços MAC de origem e encaminha quadros segundo a tabela MAC. Um roteador decide o próximo salto a partir do endereço IP de destino e da tabela de rotas. Firewall pode operar em diferentes camadas e profundidades de inspeção. A banca frequentemente troca esses papéis.

\subsection{IPv4, CIDR e subnetting}

IPv4 possui 32 bits. Em CIDR, \texttt{/n} indica o comprimento do prefixo. O número total de endereços de um bloco é:
\[
N=2^{32-n}.
\]
No uso tradicional de sub-redes, dois endereços são reservados para rede e broadcast, de modo que hosts utilizáveis são $N-2$, ressalvadas situações especiais como enlaces ponto a ponto.

\begin{exemplo}[sub-rede /27]
Para \texttt{192.0.2.64/27}, há $2^{5}=32$ endereços. O incremento é 32: blocos .0, .32, .64, .96 etc. O bloco contém .64--.95; rede = .64; broadcast = .95; hosts usuais = .65--.94.
\end{exemplo}

\subsubsection{Agregação e longest prefix match}

Roteadores preferem, entre rotas compatíveis, o prefixo mais específico. Assim, \texttt{10.0.1.0/24} prevalece sobre \texttt{10.0.0.0/8} para destino 10.0.1.20. CIDR permite sumarização de rotas quando blocos são contíguos e alinhados, reduzindo tabelas e propagação de rotas.

\subsection{IPv6}

IPv6 possui 128 bits e elimina broadcast tradicional, recorrendo a multicast e Neighbor Discovery. Endereços link-local são importantes para operações locais e autoconfiguração. A representação hexadecimal permite omitir zeros à esquerda e comprimir uma única sequência contínua de grupos zero com \texttt{::}. IPv6 não é simplesmente ``IPv4 com mais endereços'': muda mecanismos de descoberta, autoconfiguração e tratamento de fragmentação.

\subsection{ARP, ICMP, DHCP e DNS}

\termo{ARP} resolve IPv4 conhecido para endereço MAC no enlace local. Em IPv6, Neighbor Discovery desempenha funções correspondentes usando ICMPv6. \termo{ICMP} transporta mensagens de controle e diagnóstico; bloquear indiscriminadamente ICMP pode quebrar funcionalidades como Path MTU Discovery.

\termo{DHCP} automatiza parâmetros de rede por concessões. Em IPv4, a sequência didática DORA — Discover, Offer, Request, Acknowledge — ajuda a entender o processo, mas é necessário conhecer também relay para atender clientes fora do domínio de broadcast do servidor.

\termo{DNS} é uma base distribuída hierárquica. Resolvedores recursivos consultam servidores autoritativos e usam cache. Registros importantes: A, AAAA, CNAME, MX, NS, PTR e TXT. TTL controla tempo de cache, não segurança. DNSSEC adiciona validação criptográfica de autenticidade/integridade dos dados DNS, não confidencialidade das consultas.

\section{TCP, UDP e comportamento de transporte}

\subsection{TCP}

TCP é orientado à conexão e fornece fluxo confiável de bytes. O three-way handshake sincroniza números iniciais de sequência. ACKs, temporizadores e retransmissões tratam perdas; números de sequência restauram ordem; janela do receptor implementa controle de fluxo; mecanismos de congestionamento ajustam envio segundo condições da rede.

O throughput de conexão de longa distância não depende apenas de banda nominal. Produto largura de banda-atraso, RTT, perdas, tamanho de janela e congestionamento podem limitar desempenho. Em auditoria de desempenho, uma média de utilização de link baixa não prova ausência de gargalo em conexões sensíveis a RTT ou picos.

\subsection{UDP}

UDP é datagrama não orientado à conexão, sem confiabilidade, ordenação ou controle de congestionamento fornecidos pelo próprio protocolo. Isso reduz overhead e deixa à aplicação decisões de confiabilidade. DNS tradicional, voz/streaming em muitos cenários e QUIC utilizam UDP. QUIC demonstra por que ``UDP = aplicação não confiável'' é afirmação errada: a aplicação/protocolo superior pode adicionar mecanismos próprios.

\begin{cebraspe}
Evite frases absolutas como ``TCP é lento'' e ``UDP é rápido''. A diferença estrutural está nos serviços oferecidos. O desempenho real depende de protocolo superior, implementação, perdas, RTT, criptografia, congestionamento e padrão de tráfego.
\end{cebraspe}

\section{LAN, WAN, switching e routing}

\subsection{Ethernet, domínio de colisão e broadcast}

Em Ethernet comutada, cada porta opera como domínio de colisão separado em full duplex, enquanto o domínio de broadcast se estende pela VLAN. Switch aprende MAC pela origem; destino desconhecido é inundado no domínio apropriado. Loops físicos podem causar broadcast storm, duplicação de quadros e instabilidade de MAC.

\subsection{STP/RSTP e agregação}

STP constrói uma topologia lógica sem loop, elegendo root bridge e bloqueando caminhos redundantes. RSTP melhora convergência. LACP agrega enlaces físicos em um canal lógico, aumentando capacidade e redundância conforme algoritmo de distribuição; não implica que um único fluxo use automaticamente a soma integral de todos os links.

\subsection{VLAN e 802.1Q}

VLAN cria segmentação lógica de camada 2. Porta access pertence tipicamente a uma VLAN; trunk transporta múltiplas VLANs com tags 802.1Q. Comunicação entre VLANs exige roteamento. Para segurança, VLAN precisa ser acompanhada de políticas de camada 3/4; regra permissiva \texttt{any-any} entre VLANs anula grande parte do benefício de contenção.

\subsection{Roteamento}

Roteamento estático é simples e previsível, mas pouco adaptativo em topologias extensas. Protocolos dinâmicos trocam informação e calculam caminhos. O edital fala em noções de routing; portanto, compreenda diferença entre IGP e EGP, métricas, convergência, rota default e princípio de longest prefix match. OSPF é IGP de estado de enlace; BGP é protocolo interdomínio baseado em políticas e atributos, não simplesmente ``o protocolo de menor caminho''.

\subsection{MPLS}

MPLS insere rótulos para encaminhamento em redes de provedores, permitindo engenharia de tráfego e serviços como VPNs. Ele se posiciona conceitualmente entre mecanismos de enlace e rede. MPLS VPN não deve ser confundida automaticamente com VPN criptográfica: isolamento lógico no provedor não implica criptografia de conteúdo.

\section{Redes sem fio e controle de acesso}

802.11 define WLAN. O meio é compartilhado e sujeito a interferência, atenuação, roaming e contenção. Segurança evoluiu de WEP — quebrado e inadequado — para WPA/WPA2 e padrões posteriores. O edital nomeia WEP, WPA e WPA2, portanto domine a evolução e as fragilidades de WEP/TKIP frente a AES/CCMP.

802.1X implementa controle de acesso baseado em porta. Há três papéis: \termo{supplicant} (cliente), \termo{authenticator} (switch/AP) e \termo{authentication server} (frequentemente RADIUS). EAP é uma estrutura que transporta métodos de autenticação; não é um único método específico.

\section{Gerenciamento de redes: SMI, SNMP e MIB}

Gerenciamento de redes envolve coleta de estado, configuração, falhas, desempenho, segurança e inventário. Em SNMP, um manager consulta agentes que expõem objetos organizados em MIB. A SMI define convenções de nomes, tipos e estrutura desses objetos.

Operações clássicas incluem GET, GETNEXT/GETBULK, SET e notificações TRAP/INFORM. Polling frequente aumenta visibilidade, mas também carga; traps são assíncronos, porém podem ser perdidos conforme versão/transporte. SNMPv1/v2c dependem de community strings com proteção insuficiente; SNMPv3 pode fornecer autenticação e privacidade conforme configuração.

\begin{exemplo}[auditoria de monitoramento]
Não basta verificar que ``há Zabbix/PRTG/SNMP''. O auditor deve examinar cobertura: quais ativos estão monitorados, quais métricas, thresholds, retenção, sincronização de tempo, tratamento de alertas e evidência de resposta. Ferramenta instalada sem processo não constitui controle efetivo.
\end{exemplo}

\section{Segurança em redes}

O edital repete parte de segurança de redes dentro de infraestrutura. Firewalls aplicam regras; IDS detecta; IPS pode bloquear inline; proxy intermedeia; NAT traduz endereços; VPN cria canal lógico protegido conforme tecnologia. Nenhum deles substitui os demais.

Ataques cobrados: spoofing, flood, DoS/DDoS, phishing e malwares. Spoofing falsifica identidade/origem em determinado protocolo; flood esgota recursos por volume; DDoS distribui a origem. Segmentação, rate limiting, anti-DDoS, autenticação forte, filtragem e observabilidade são controles complementares.

\section{Criptografia aplicada à infraestrutura}

Criptografia simétrica usa segredo compartilhado e é eficiente para grandes volumes; assimétrica usa par de chaves e habilita autenticação, troca de chaves e assinatura. Sistemas reais combinam ambas. Hash criptográfico produz digest e não é cifra reversível.

Assinatura digital vincula mensagem/chave privada para apoiar autenticidade e integridade; não cifra necessariamente o conteúdo. Certificação digital vincula identidade a chave pública por cadeia de confiança. TLS protege dados em trânsito e utiliza mecanismos híbridos; VPNs IPsec podem proteger tráfego em camada de rede.

\section{Sistemas operacionais}

\subsection{Processos, threads e escalonamento}

Processo é uma instância de programa com contexto de execução, espaço de endereçamento e recursos. Threads de um mesmo processo compartilham memória e recursos, mas possuem pilha, registradores e contador de programa próprios. Concorrência cria necessidade de sincronização.

Estados clássicos: novo, pronto, executando, bloqueado/esperando e terminado. O escalonador seleciona processos prontos. Algoritmos como FCFS, SJF/SRTF, Round Robin e prioridades apresentam trade-offs de tempo de resposta, throughput, justiça e starvation.

\subsection{Condição de corrida, exclusão mútua e deadlock}

Condição de corrida ocorre quando o resultado depende da ordem/interleaving de acessos concorrentes. Mutex, semáforos e monitores coordenam seções críticas. Deadlock exige, no modelo clássico, quatro condições simultâneas: exclusão mútua, posse e espera, ausência de preempção e espera circular. Prevenção rompe pelo menos uma; detecção permite ocorrer e identifica ciclos; avoidance depende de conhecimento de demandas e estado seguro.

\subsection{Memória virtual}

Paginação divide memória virtual em páginas e física em frames. Tabela de páginas traduz endereços; TLB armazena traduções recentes. Page fault ocorre quando a página necessária não está residente e deve ser carregada. Algoritmos de substituição buscam explorar localidade; LRU é referência conceitual, embora implementações reais usem aproximações.

\termo{Thrashing} ocorre quando o working set excede memória disponível e o sistema passa grande parte do tempo em page faults. Adicionar CPU não resolve necessariamente; reduzir multiprogramação ou aumentar memória pode ser mais efetivo.

\subsection{Entrada/saída e sistemas de arquivos}

I/O envolve drivers, interrupções, buffers, filas e DMA em várias arquiteturas. Cache e write-back melhoram desempenho, mas exigem políticas de persistência. Sistemas de arquivos organizam arquivos, diretórios, metadados e permissões; journaling melhora recuperação de metadados após falha, mas não substitui backup.

\subsection{Linux}

Domine usuários/grupos, permissões rwx, octal, umask, ACLs, sudo, processos, systemd, logs, pacotes, filesystem, rede e SSH. Permissões tradicionais são apenas uma camada: capacidades, SELinux/AppArmor e ACLs podem alterar o resultado efetivo. Em auditoria, \texttt{chmod 777} é indício de privilégio excessivo, não solução adequada de acesso.

\subsection{Windows Server}

Estude serviços, Event Viewer, NTFS, compartilhamentos, PowerShell, atualização, firewall, RDP, Active Directory, GPO e autenticação. Permissão efetiva pode resultar da combinação de ACL NTFS e permissão de compartilhamento; ao acessar pela rede, aplica-se a combinação mais restritiva entre as camadas pertinentes.

\section{Active Directory, LDAP e interoperabilidade}

LDAP é protocolo de diretório; AD é serviço de diretório que integra LDAP, Kerberos, DNS e outros componentes. Domínios formam limites administrativos lógicos; árvores e florestas organizam namespaces e relações de confiança. O Global Catalog acelera consultas a objetos na floresta.

Kerberos usa KDC, TGT e service tickets. Sincronização de relógio é relevante para reduzir replay. GPO centraliza configuração. Interoperabilidade pode envolver LDAP, Kerberos, SAML/OIDC em camadas superiores e sincronização/federação de identidades.

\begin{cebraspe}
LDAP não é ``o banco do Active Directory'' nem sinônimo de AD. Kerberos não é protocolo de diretório. DNS não é opcional em um domínio AD moderno: é infraestrutura crítica para localização de serviços.
\end{cebraspe}

\section{Virtualização e cloud computing}

Hipervisor tipo 1 executa diretamente sobre hardware; tipo 2 sobre sistema hospedeiro, na classificação tradicional. Virtualização melhora consolidação e isolamento, mas amplia risco de concentração: falha de host, storage ou management plane pode afetar muitas VMs.

Snapshots preservam estado em ponto no tempo e são úteis para rollback/backup auxiliar, mas longas cadeias degradam desempenho e snapshots no mesmo storage não são cópia independente. Live migration move workloads com interrupção mínima conforme plataforma, exigindo rede, storage ou mecanismos equivalentes adequados.

Cloud computing adiciona elasticidade, autosserviço, pooling e medição de uso. O aprofundamento de IaaS/PaaS/SaaS e arquitetura cloud está no capítulo específico de nuvem.

\section{Servidores de aplicação, HA e escalabilidade}

\subsection{Topologia em camadas}

Ambiente típico separa entrada/reverse proxy, camada de aplicação, serviços de estado/cache/mensageria e banco. Alta disponibilidade exige identificar dependências em cada camada. Dois servidores web atrás de balanceador não tornam o serviço altamente disponível se banco, DNS ou storage forem únicos.

\subsection{Balanceamento}

Balanceadores distribuem requisições por round-robin, least connections, hashing ou políticas. Health checks retiram instâncias inadequadas. Sticky sessions mantêm afinidade, mas podem criar desbalanceamento e complicar failover; externalizar estado costuma aumentar resiliência.

\subsection{Failover e replicação}

Failover transfere serviço para redundância. Replicação síncrona prioriza consistência com custo de latência/disponibilidade em falhas de comunicação; assíncrona reduz latência do primário, mas pode aumentar RPO. A escolha deve refletir requisitos de negócio.

\subsection{Desempenho}

Métricas importantes: latência média e percentis, throughput, erro, saturação, filas, CPU, memória, I/O e conexões. Percentil 95/99 revela cauda que a média esconde. Little, em sistema estável, relaciona itens médios no sistema, taxa e tempo: $L=\lambda W$.

\section{Data center e ambientes de missão crítica}

Data center é sistema de dependências: energia, climatização, incêndio, conectividade, segurança física, racks, cabeamento, compute, storage e gestão. Redundância deve ser analisada por \termo{domínio de falha}. Duas fontes ligadas ao mesmo UPS ou dois links passando pelo mesmo duto continuam compartilhando causa comum.

Conceitos de topologia N, N+1, 2N e equivalentes expressam níveis de capacidade/redundância, mas a disponibilidade real depende de arquitetura, manutenção e operação. Gestão de mudança e inventário físico são controles tão importantes quanto equipamento redundante.

\section{Armazenamento: discos, interfaces, RAID, NAS e SAN}

\subsection{Disco e interface}

HDD usa mídia magnética e partes mecânicas, favorecendo capacidade/custo; SSD usa flash, reduz latência e aumenta IOPS, mas possui características de desgaste e controladora. SATA, SAS e NVMe representam interfaces/protocolos com perfis distintos. IOPS, throughput e latência devem ser analisados conforme tamanho e padrão de I/O.

\subsection{RAID}

\begin{table}[ht]
\centering
\small
\begin{tabularx}{\textwidth}{l c c X}
\toprule
RAID & Mínimo & Tolerância típica & Característica \\
\midrule
0 & 2 & nenhuma & striping; desempenho/capacidade, sem redundância. \\
1 & 2 & um disco por espelho & espelhamento; leitura pode ganhar desempenho. \\
5 & 3 & 1 disco & paridade distribuída; penalidade de escrita. \\
6 & 4 & 2 discos & dupla paridade; maior proteção a falhas durante rebuild. \\
10 & 4 & depende dos pares & espelhos com striping; bom I/O, custo de 50\% da capacidade bruta. \\
\bottomrule
\end{tabularx}
\end{table}

RAID reduz risco de falha de disco conforme nível, mas não cobre exclusão, corrupção lógica, ransomware, desastre, erro administrativo ou comprometimento de storage. Portanto, \textbf{RAID não é backup}.

\subsection{NAS e SAN}

NAS oferece arquivo via SMB/NFS e apresenta namespace de arquivos ao cliente. SAN oferece bloco por tecnologias como Fibre Channel ou iSCSI e apresenta LUNs/volumes ao host. ``SAN é mais rápida'' não é regra universal: desempenho depende de mídia, protocolo, rede, controladora, cache, workload e desenho.

\section{Backup, restauração e deduplicação}

\subsection{Full, incremental e diferencial}

Backup full copia conjunto definido. Incremental copia alterações desde o último backup de qualquer tipo; restauração exige full + cadeia de incrementais. Diferencial copia alterações desde o último full; cresce ao longo do ciclo, mas restauração costuma exigir full + último diferencial.

\subsection{RPO e RTO}

\termo{RPO} é perda temporal tolerável de dados; \termo{RTO} é tempo-alvo de restabelecimento. Um RPO de 15 min é incompatível, sem mecanismos adicionais, com único ponto de recuperação diário. RTO depende de capacidade de reconstruir infraestrutura, recuperar dados, validar consistência e retornar tráfego.

\subsection{3-2-1 e imutabilidade}

A regra 3-2-1 é heurística: três cópias, em dois tipos de mídia/local, uma fora do ambiente principal. Estratégias modernas adicionam cópia offline/imutável e verificação. O ponto central é diversidade de domínio de falha e proteção contra credenciais comprometidas.

\subsection{Teste de restauração}

Job ``sucesso'' apenas prova que software terminou uma tarefa, não que dados são recuperáveis. Testes devem validar catálogo, cadeia, chaves de criptografia, consistência de aplicação e tempo real de recuperação. Para sistemas críticos, teste de DR deve incluir dependências e procedimentos de retorno.

Deduplicação elimina blocos redundantes e reduz armazenamento/tráfego, mas cria dependência de metadados e não substitui múltiplas cópias independentes.

\section{Contêinerização e DevOps}

Contêineres usam namespaces/cgroups e outros mecanismos para isolamento e controle de recursos, compartilhando kernel do host. Imagem é artefato em camadas; contêiner é instância executável. Imagens devem ser versionadas, mínimas, escaneadas e produzidas por pipeline confiável.

Riscos recorrentes: tag mutável \texttt{latest}, execução privilegiada, montagem do socket Docker, segredo em imagem, base desatualizada e registry sem controle. Cadeia de suprimentos exige origem, assinatura/proveniência, SBOM quando aplicável e promoção controlada.

DevOps integra desenvolvimento e operação por automação, feedback e responsabilidade compartilhada. Para Auditor/TI, observe CI/CD, segregação, credenciais, trilha, rollback, testes, monitoramento e aprovação baseada em risco.

\section{Mapa mental}

\mapamental{Infraestrutura}{
 child { node[concept] {Redes} child {node[concept]{TCP/IP}} child {node[concept]{VLAN/routing}} }
 child { node[concept] {SO} child {node[concept]{processos}} child {node[concept]{memória/I-O}} }
 child { node[concept] {Identidade} child {node[concept]{AD/LDAP}} child {node[concept]{Kerberos/GPO}} }
 child { node[concept] {Missão crítica} child {node[concept]{HA}} child {node[concept]{data center}} }
 child { node[concept] {Storage} child {node[concept]{RAID/NAS/SAN}} child {node[concept]{backup/RPO}} }
 child { node[concept] {Plataforma} child {node[concept]{virtualização}} child {node[concept]{containers/DevOps}} }
}

\begin{resumo}
\textbf{Fixação:} não confunda redundância com backup; VLAN com política de segurança; disponibilidade com durabilidade; LDAP com Active Directory; snapshot com cópia independente; média de desempenho com comportamento de cauda. Em auditoria, identifique sempre o domínio de falha, o RPO/RTO exigido e a evidência de que o controle funciona de fato.
\end{resumo}

\section{Banco de questões — Infraestrutura de TI}

\subsection*{N1 — Fundamento competitivo}

\begin{questao}{\nivel{N1} 10.01 — Cebraspe/Certo ou Errado}
Em uma rede IPv4 \texttt{/27}, o bloco contém 32 endereços; em uso convencional, 30 podem ser atribuídos a hosts, pois um endereço identifica a rede e outro o broadcast.
\end{questao}

\begin{questao}{\nivel{N1} 10.02 — Cebraspe/Certo ou Errado}
O fato de o UDP não fornecer retransmissão e ordenação nativamente implica que qualquer aplicação baseada em UDP seja, necessariamente, incapaz de oferecer entrega confiável.
\end{questao}

\begin{questao}{\nivel{N1} 10.03 — Cebraspe/Certo ou Errado}
Uma VLAN separa domínios de broadcast em camada 2, mas não impede, por si só, comunicação entre redes quando existe roteamento inter-VLAN com política permissiva.
\end{questao}

\begin{questao}{\nivel{N1} 10.04 — Cebraspe/Certo ou Errado}
Em um switch Ethernet, o endereço MAC de origem recebido em um quadro é usado para aprendizado da tabela de encaminhamento, ao passo que o endereço MAC de destino é usado na decisão de encaminhamento.
\end{questao}

\begin{questao}{\nivel{N1} 10.05 — Cebraspe/Certo ou Errado}
SNMPv3 pode fornecer autenticação e privacidade, mas sua mera adoção não garante gerenciamento seguro se credenciais, perfis e acesso à rede de gerenciamento forem mal configurados.
\end{questao}

\begin{questao}{\nivel{N1} 10.06 — Cebraspe/Certo ou Errado}
LDAP e Active Directory são termos equivalentes, pois ambos designam um serviço de diretório completo e independente de outros protocolos.
\end{questao}

\begin{questao}{\nivel{N1} 10.07 — Cebraspe/Certo ou Errado}
RAID 6 tolera, em condições normais do arranjo, a falha de até dois discos, mas isso não o transforma em substituto de backup.
\end{questao}

\begin{questao}{\nivel{N1} 10.08 — Cebraspe/Certo ou Errado}
Uma snapshot armazenada no mesmo storage da máquina virtual constitui cópia independente suficiente contra perda total do storage original.
\end{questao}

\begin{questao}{\nivel{N1} 10.09 — Cebraspe/Certo ou Errado}
Em um sistema paginado, thrashing caracteriza situação em que o tempo gasto com faltas e substituição de páginas cresce a ponto de reduzir significativamente o trabalho útil da CPU.
\end{questao}

\begin{questao}{\nivel{N1} 10.10 — Cebraspe/Certo ou Errado}
Duas máquinas virtuais executadas no mesmo host físico eliminam o ponto único de falha do hardware, pois cada VM constitui um domínio de falha independente.
\end{questao}

\subsection*{N2 — Aplicação}

\begin{questao}{\nivel{N2} 10.11 — FCC/subnetting}
Um segmento precisa suportar 50 hosts IPv4 em uma única sub-rede, sem considerar endereçamento especial além de rede e broadcast. O menor prefixo CIDR adequado é:
\begin{enumerate}[label=\Alph*)]
\item /28
\item /27
\item /26
\item /25
\item /24
\end{enumerate}
\end{questao}

\begin{questao}{\nivel{N2} 10.12 — FGV/roteamento}
A tabela de rotas de um roteador contém \texttt{10.0.0.0/8}, \texttt{10.20.0.0/16} e \texttt{10.20.30.0/24}. Um pacote destinado a \texttt{10.20.30.55} será encaminhado, em regra, pela rota:
\begin{enumerate}[label=\Alph*)]
\item \texttt{10.0.0.0/8}, por ser a mais antiga;
\item \texttt{10.20.0.0/16}, por possuir métrica intermediária;
\item \texttt{10.20.30.0/24}, por ser o prefixo mais específico;
\item default, porque existem rotas sobrepostas;
\item aleatoriamente entre as três.
\end{enumerate}
\end{questao}

\begin{questao}{\nivel{N2} 10.13 — Cebraspe/Certo ou Errado}
Um trunk 802.1Q permite transportar múltiplas VLANs entre switches, mas não substitui o roteamento necessário para comunicação entre VLANs distintas.
\end{questao}

\begin{questao}{\nivel{N2} 10.14 — FGV/802.1X}
Uma instituição deseja liberar portas cabeadas e acesso Wi-Fi apenas após autenticação centralizada. A arquitetura mais aderente é:
\begin{enumerate}[label=\Alph*)]
\item 802.1X com supplicant, autenticador e backend RADIUS/EAP;
\item STP com DNS recursivo;
\item SNMP com community string pública;
\item LACP com NAT;
\item MPLS com DHCP relay.
\end{enumerate}
\end{questao}

\begin{questao}{\nivel{N2} 10.15 — Cebraspe/Certo ou Errado}
O bloqueio indiscriminado de ICMP pode prejudicar funcionalidades legítimas de rede, como diagnóstico e mecanismos relacionados a descoberta de MTU de caminho.
\end{questao}

\begin{questao}{\nivel{N2} 10.16 — FCC/processos}
Um processo em execução solicita leitura de disco e não pode prosseguir até a conclusão do I/O. O estado mais coerente, imediatamente após a solicitação, é:
\begin{enumerate}[label=\Alph*)]
\item pronto;
\item bloqueado/esperando;
\item terminado;
\item zombie;
\item executando com prioridade máxima.
\end{enumerate}
\end{questao}

\begin{questao}{\nivel{N2} 10.17 — FGV/deadlock}
Dois processos mantêm recursos distintos e cada um espera indefinidamente pelo recurso mantido pelo outro. O cenário evidencia diretamente:
\begin{enumerate}[label=\Alph*)]
\item thrashing;
\item starvation sem espera circular;
\item uma cadeia de espera circular compatível com deadlock;
\item fragmentação externa;
\item falha de paginação.
\end{enumerate}
\end{questao}

\begin{questao}{\nivel{N2} 10.18 — Cebraspe/Certo ou Errado}
Em Linux, conceder \texttt{chmod 777} para resolver um erro de acesso pode eliminar a causa imediata, mas cria risco de privilégio excessivo e não deve ser tratado como solução segura por padrão.
\end{questao}

\begin{questao}{\nivel{N2} 10.19 — FGV/AD}
Em um domínio Active Directory, usuários conseguem autenticar-se localmente, mas falham ao localizar controladores de domínio após alteração incorreta do DNS. Isso é coerente porque:
\begin{enumerate}[label=\Alph*)]
\item o AD depende de DNS para localização de serviços e controladores;
\item LDAP substitui integralmente o DNS;
\item Kerberos não utiliza nomes de serviço;
\item GPO dispensa resolução de nomes;
\item o Global Catalog funciona sem qualquer infraestrutura DNS.
\end{enumerate}
\end{questao}

\begin{questao}{\nivel{N2} 10.20 — FCC/RAID}
Um arranjo RAID 5 utiliza quatro discos de 6 TB. Desprezando overhead e diferenças de unidade, a capacidade útil aproximada é:
\begin{enumerate}[label=\Alph*)]
\item 6 TB
\item 12 TB
\item 18 TB
\item 24 TB
\item 30 TB
\end{enumerate}
\end{questao}

\begin{questao}{\nivel{N2} 10.21 — Cebraspe/Certo ou Errado}
Em RAID 10 com quatro discos, a tolerância a duas falhas simultâneas depende de quais discos falham; perder os dois membros do mesmo espelho pode causar perda do conjunto.
\end{questao}

\begin{questao}{\nivel{N2} 10.22 — FGV/NAS e SAN}
Uma aplicação exige volume em bloco apresentado ao servidor como LUN, enquanto outra precisa de compartilhamento de arquivos SMB para usuários. As tecnologias mais aderentes são, respectivamente:
\begin{enumerate}[label=\Alph*)]
\item NAS e SAN;
\item SAN e NAS;
\item CDN e SAN;
\item SAN e DNS;
\item NAS e DHCP.
\end{enumerate}
\end{questao}

\begin{questao}{\nivel{N2} 10.23 — FCC/backup}
Uma estratégia utiliza backup full aos domingos e incrementais de segunda a sábado. Para restaurar o estado de sexta-feira, normalmente será necessário:
\begin{enumerate}[label=\Alph*)]
\item apenas o incremental de sexta;
\item o full de domingo e todos os incrementais de segunda a sexta;
\item o full de domingo e somente o incremental de segunda;
\item todos os fulls do mês;
\item apenas o último snapshot de memória.
\end{enumerate}
\end{questao}

\begin{questao}{\nivel{N2} 10.24 — Cebraspe/Certo ou Errado}
Um RPO de 15 minutos pode ser incompatível com a realização de apenas um backup completo diário, na ausência de mecanismos adicionais de replicação ou captura contínua de alterações.
\end{questao}

\begin{questao}{\nivel{N2} 10.25 — FGV/containers}
Uma aplicação mantém estado de sessão exclusivamente na memória local de cada contêiner e passa a ser escalada horizontalmente. A medida arquitetural mais adequada é:
\begin{enumerate}[label=\Alph*)]
\item fixar o tráfego permanentemente a um único contêiner;
\item externalizar/replicar o estado ou adotar desenho stateless compatível;
\item desabilitar health checks;
\item substituir TCP por UDP;
\item armazenar sessão em camada temporária da imagem.
\end{enumerate}
\end{questao}

\subsection*{N3 — Integração de concurso}

\begin{questao}{\nivel{N3} 10.26 — FGV/assertivas de redes}
Considere:

I. STP/RSTP busca impedir loops de camada 2.

II. LACP pode agregar enlaces físicos em um canal lógico.

III. A existência de LACP implica que um único fluxo TCP utilizará sempre a soma integral da largura de banda de todos os membros.

Assinale a alternativa correta.
\begin{enumerate}[label=\Alph*)]
\item Apenas I.
\item Apenas II.
\item Apenas I e II.
\item Apenas II e III.
\item I, II e III.
\end{enumerate}
\end{questao}

\begin{questao}{\nivel{N3} 10.27 — Cebraspe/Certo ou Errado}
Uma rede segmentada em VLANs continua vulnerável a movimento lateral se o firewall ou roteador inter-VLAN permitir tráfego amplo entre as zonas, pois a segmentação lógica sem política restritiva não garante isolamento efetivo.
\end{questao}

\begin{questao}{\nivel{N3} 10.28 — FCC/desempenho TCP}
Um link de 1 Gbit/s entre dois datacenters possui RTT elevado e baixa perda. Uma única transferência TCP atinge apenas fração da banda disponível. A explicação mais adequada é:
\begin{enumerate}[label=\Alph*)]
\item banda nominal é o único fator relevante, logo o comportamento é impossível;
\item janela TCP, RTT, congestionamento e produto largura de banda-atraso podem limitar o throughput de um fluxo;
\item UDP resolveria automaticamente qualquer problema de confiabilidade;
\item VLAN reduz throughput de longa distância por definição;
\item DNS é o principal limitador de transferência contínua após a conexão.
\end{enumerate}
\end{questao}

\begin{questao}{\nivel{N3} 10.29 — Cebraspe/Certo ou Errado}
O uso de MPLS VPN por uma operadora fornece isolamento lógico de tráfego, mas não deve ser confundido automaticamente com criptografia fim a fim do conteúdo transportado.
\end{questao}

\begin{questao}{\nivel{N3} 10.30 — FGV/SNMP}
Uma rede de gerenciamento usa SNMPv3, mas todos os equipamentos compartilham o mesmo usuário, a mesma chave e aceitam acesso de qualquer VLAN. A melhor avaliação é:
\begin{enumerate}[label=\Alph*)]
\item SNMPv3 torna o desenho seguro independentemente da configuração;
\item o protocolo oferece controles melhores, mas a configuração ainda viola segregação e mínimo privilégio;
\item deve-se voltar a SNMPv2c para reduzir risco;
\item MIB e SMI substituem autenticação;
\item community strings tornam-se mais seguras em SNMPv3.
\end{enumerate}
\end{questao}

\begin{questao}{\nivel{N3} 10.31 — FCC/memória virtual}
Um servidor apresenta alta utilização de disco, taxa elevada de page faults e CPU frequentemente ociosa esperando I/O. Aumentar apenas a frequência da CPU tende a produzir pouco efeito porque o gargalo principal pode ser:
\begin{enumerate}[label=\Alph*)]
\item excesso de memória disponível;
\item thrashing por pressão de memória/working set;
\item tabela MAC cheia;
\item DNS recursivo;
\item RAID 1.
\end{enumerate}
\end{questao}

\begin{questao}{\nivel{N3} 10.32 — Cebraspe/Certo ou Errado}
MV e contêiner oferecem isolamento por mecanismos distintos; em geral, uma VM inclui sistema operacional convidado sobre hipervisor, enquanto contêineres compartilham o kernel do host.
\end{questao}

\begin{questao}{\nivel{N3} 10.33 — FGV/HA}
Dois servidores de aplicação estão em hosts físicos distintos, mas ambos dependem do mesmo storage sem controladora redundante. A arquitetura possui:
\begin{enumerate}[label=\Alph*)]
\item alta disponibilidade completa, pois o compute é redundante;
\item ponto único de falha no storage compartilhado;
\item redundância 2N em todas as camadas;
\item tolerância a desastre regional;
\item independência total entre os servidores.
\end{enumerate}
\end{questao}

\begin{questao}{\nivel{N3} 10.34 — Cebraspe/Certo ou Errado}
A disponibilidade de um serviço não pode ser inferida apenas pela redundância de servidores, porque dependências comuns de DNS, storage, banco, energia ou telecomunicações podem permanecer como pontos únicos de falha.
\end{questao}

\begin{questao}{\nivel{N3} 10.35 — FCC/RPO-RTO}
Um sistema possui RPO de 30 minutos e RTO de 1 hora. O restore test mais recente recuperou dados de 10 minutos antes da falha, mas o serviço demorou 3 horas para voltar. Nesse teste:
\begin{enumerate}[label=\Alph*)]
\item RPO e RTO foram cumpridos;
\item somente o RPO foi cumprido;
\item somente o RTO foi cumprido;
\item nenhum foi cumprido;
\item não é possível avaliar nenhum dos dois.
\end{enumerate}
\end{questao}

\begin{questao}{\nivel{N3} 10.36 — FGV/backup}
Um banco usa RAID 6, snapshots locais a cada hora e backup externo semanal. A administração afirma que o risco de ransomware está totalmente coberto. A melhor avaliação é:
\begin{enumerate}[label=\Alph*)]
\item correta, porque RAID 6 protege contra ransomware;
\item correta, porque snapshots locais são imunes à exclusão;
\item inadequada; é preciso avaliar independência, imutabilidade, credenciais e janela de perda do backup externo;
\item correta se o banco for ACID;
\item correta se o storage usar SSD.
\end{enumerate}
\end{questao}

\begin{questao}{\nivel{N3} 10.37 — Cebraspe/Certo ou Errado}
A deduplicação pode reduzir consumo de armazenamento, mas aumenta dependência de metadados e não substitui a existência de cópias em domínios de falha independentes.
\end{questao}

\begin{questao}{\nivel{N3} 10.38 — FGV/Kerberos}
Uma conta administrativa utiliza a mesma identidade para e-mail, navegação e administração de controladores de domínio. Ainda que Kerberos funcione corretamente, o desenho é inadequado porque:
\begin{enumerate}[label=\Alph*)]
\item a mesma conta aumenta superfície de exposição de credencial privilegiada;
\item Kerberos proíbe e-mail;
\item LDAP deixa de funcionar;
\item GPO exige uma conta por aplicação;
\item DNS não suporta contas administrativas.
\end{enumerate}
\end{questao}

\begin{questao}{\nivel{N3} 10.39 — Cebraspe/Certo ou Errado}
Uma liveness probe que dependa de serviço externo pode provocar reinícios desnecessários de instâncias saudáveis durante indisponibilidade da dependência, amplificando a falha.
\end{questao}

\begin{questao}{\nivel{N3} 10.40 — FCC/capacidade}
A CPU média de um servidor é 22\%, mas o p99 de CPU chega a 100\% durante fechamento mensal, período em que o SLA é violado. A conclusão mais adequada é:
\begin{enumerate}[label=\Alph*)]
\item a média prova excesso de capacidade;
\item planejamento deve considerar picos, filas, percentis e janela crítica, não apenas média global;
\item CPU não influencia SLA;
\item o servidor está subutilizado em qualquer momento;
\item a única solução possível é scale-up.
\end{enumerate}
\end{questao}

\subsection*{N4 — Auditoria / avançado}

\begin{questao}{\nivel{N4} 10.41 — Cebraspe/caso integrado}
Em auditoria de continuidade, a equipe encontra: RAID 6; snapshots no mesmo storage; backup externo mensal; RPO declarado de 15 minutos; e nenhum restore test nos últimos 18 meses. Julgue o item:

A presença de RAID 6 e snapshots frequentes é suficiente para demonstrar o atendimento ao RPO, ainda que o storage seja perdido integralmente.
\end{questao}

\begin{questao}{\nivel{N4} 10.42 — FGV/rede corporativa}
Usuários, servidores e administração estão em VLANs distintas, porém a política do firewall interno é \texttt{any-any} entre todas as zonas. Em um teste de comprometimento de estação de usuário, o atacante alcança RDP e SMB dos servidores. O achado principal é:
\begin{enumerate}[label=\Alph*)]
\item inexistência de VLANs;
\item segmentação nominal sem política efetiva de contenção, permitindo movimento lateral;
\item falha exclusiva de DNS;
\item ausência de STP;
\item necessidade de trocar Ethernet por MPLS.
\end{enumerate}
\end{questao}

\begin{questao}{\nivel{N4} 10.43 — FCC/AD e logs}
Administradores de domínio usam contas compartilhadas e conseguem excluir logs locais dos controladores. Qual conjunto de controles é mais adequado?
\begin{enumerate}[label=\Alph*)]
\item contas nominativas, separação de privilégio, MFA/PAM e envio de logs para repositório protegido;
\item aumentar apenas o tamanho das senhas compartilhadas;
\item manter contas compartilhadas e reduzir retenção de logs;
\item desabilitar auditoria para melhorar desempenho;
\item migrar os controladores para RAID 0.
\end{enumerate}
\end{questao}

\begin{questao}{\nivel{N4} 10.44 — Cebraspe/Certo ou Errado}
Dois datacenters podem continuar compartilhando uma causa comum de falha se ambos dependem do mesmo provedor e do mesmo ponto físico de entrada de fibra, de modo que redundância geográfica aparente não garante diversidade de conectividade.
\end{questao}

\begin{questao}{\nivel{N4} 10.45 — FGV/containers}
Um cluster de contêineres aceita imagens sem assinatura, usa tag \texttt{latest} em produção e permite execução privilegiada por padrão. O risco combinado mais relevante é:
\begin{enumerate}[label=\Alph*)]
\item supply chain não verificável, baixa reprodutibilidade e privilégio excessivo;
\item apenas falta de espaço em disco;
\item somente indisponibilidade de DNS;
\item ausência de VLAN;
\item perda de capacidade RAID.
\end{enumerate}
\end{questao}

\begin{questao}{\nivel{N4} 10.46 — Cebraspe/Certo ou Errado}
Se todos os logs de infraestrutura permanecem apenas nos próprios servidores e podem ser apagados pelos mesmos administradores cujas ações são auditadas, a trilha de auditoria apresenta risco de integridade e accountability.
\end{questao}

\begin{questao}{\nivel{N4} 10.47 — FCC/HA}
Um serviço está em duas VMs no mesmo host, com banco em storage local do host e balanceador externo. A gestão classifica a arquitetura como altamente disponível. A avaliação correta é:
\begin{enumerate}[label=\Alph*)]
\item a classificação é adequada, pois existem duas VMs;
\item compute e dados compartilham o mesmo domínio de falha físico, logo a redundância é insuficiente;
\item o balanceador externo elimina qualquer SPOF interno;
\item o uso de virtualização substitui DR;
\item a única deficiência é o uso de storage local.
\end{enumerate}
\end{questao}

\begin{questao}{\nivel{N4} 10.48 — FGV/evidência de backup}
O software de backup registra ``job successful'' todas as noites. O auditor deseja concluir sobre recuperabilidade. A evidência mais apropriada é:
\begin{enumerate}[label=\Alph*)]
\item apenas o status verde do job;
\item teste documentado de restauração com validação de integridade e tempo obtido;
\item tamanho do arquivo de log;
\item capacidade bruta do storage;
\item existência de RAID no repositório.
\end{enumerate}
\end{questao}

\begin{questao}{\nivel{N4} 10.49 — Cebraspe/Certo ou Errado}
Um hash calculado sobre uma imagem forense e preservado durante a análise ajuda a demonstrar que a cópia não foi alterada, mas não prova, por si só, que a fonte original continha informação verdadeira.
\end{questao}

\begin{questao}{\nivel{N4} 10.50 — FGV/capacidade e risco}
Uma aplicação crítica possui p95 de 800 ms, p99 de 11 s e média de 320 ms. O SLA considera apenas média inferior a 1 s. Usuários relatam travamentos em operações de fechamento. A recomendação de auditoria mais adequada é:
\begin{enumerate}[label=\Alph*)]
\item manter o SLA, pois a média está cumprida;
\item revisar o indicador para incluir percentis e jornadas críticas, correlacionando latência, filas e saturação;
\item medir apenas disponibilidade;
\item remover monitoramento para reduzir overhead;
\item aumentar CPU sem diagnóstico adicional.
\end{enumerate}
\end{questao}


\section{Treino discursivo}

\begin{discursiva}[infraestrutura e evidência]
Um órgão informa RTO de 2 horas e RPO de 15 minutos para sistema crítico, mas realiza backup completo apenas uma vez por dia e não testa restauração há 18 meses. Em até 30 linhas, identifique os principais riscos, evidências adicionais que o auditor deve obter e recomendações prioritárias.
\end{discursiva}

\begin{discursiva}[alta disponibilidade]
Um ambiente possui quatro servidores de aplicação virtualizados em dois hosts, mas usa storage único e apenas um link de telecomunicações para o datacenter. Em até 30 linhas, avalie a alegação de alta disponibilidade, identifique pontos únicos de falha e proponha procedimentos de auditoria para validar o RTO declarado.
\end{discursiva}

\section{Gabarito e resoluções comentadas — Infraestrutura de TI}

\begin{solucao}{10.01}\textbf{Gabarito: CERTO.} Em /27 restam 5 bits para host: $2^5=32$ endereços totais. No uso IPv4 convencional, rede e broadcast não são atribuídos a hosts, restando 30.\end{solucao}
\begin{solucao}{10.02}\textbf{Gabarito: ERRADO.} UDP não oferece confiabilidade nativa, mas protocolos superiores podem implementá-la. QUIC é exemplo clássico: usa UDP como substrato e adiciona mecanismos próprios de entrega e congestionamento.\end{solucao}
\begin{solucao}{10.03}\textbf{Gabarito: CERTO.} VLAN separa camada 2. Se houver roteamento permissivo entre as VLANs, a comunicação continua possível. A pegadinha é confundir segmentação lógica com política de segurança.\end{solucao}
\begin{solucao}{10.04}\textbf{Gabarito: CERTO.} O switch aprende a origem e consulta o destino para decidir encaminhamento. É uma distinção clássica de prova.\end{solucao}
\begin{solucao}{10.05}\textbf{Gabarito: CERTO.} SNMPv3 melhora segurança do protocolo, mas configuração de usuários, chaves, privilégios e rede de gerenciamento continua determinante. Protocolo seguro mal configurado não produz ambiente seguro.\end{solucao}
\begin{solucao}{10.06}\textbf{Gabarito: ERRADO.} LDAP é protocolo de acesso a diretórios; Active Directory é um serviço de diretório da Microsoft que usa LDAP, Kerberos, DNS e outros componentes.\end{solucao}
\begin{solucao}{10.07}\textbf{Gabarito: CERTO.} RAID 6 usa dupla paridade e tolera duas falhas de disco dentro do desenho do arranjo, mas não protege contra exclusão, corrupção lógica, ransomware ou perda do storage inteiro.\end{solucao}
\begin{solucao}{10.08}\textbf{Gabarito: ERRADO.} Snapshot no mesmo domínio de falha pode ser perdida junto com o storage. Ela pode ajudar em rollback, mas não substitui cópia independente.\end{solucao}
\begin{solucao}{10.09}\textbf{Gabarito: CERTO.} Thrashing ocorre quando faltas de página e troca de páginas dominam o tempo, deixando pouca execução útil.\end{solucao}
\begin{solucao}{10.10}\textbf{Gabarito: ERRADO.} As VMs compartilham host, energia, hypervisor e hardware. A falha física pode derrubar ambas.\end{solucao}

\begin{solucao}{10.11}\textbf{Gabarito: C.} /26 fornece 64 endereços totais e, convencionalmente, 62 hosts utilizáveis. /27 oferece apenas 30.\end{solucao}
\begin{solucao}{10.12}\textbf{Gabarito: C.} O princípio é longest prefix match: /24 é mais específico que /16 e /8 para o destino informado.\end{solucao}
\begin{solucao}{10.13}\textbf{Gabarito: CERTO.} Trunk transporta VLANs; roteamento continua necessário para tráfego entre redes distintas.\end{solucao}
\begin{solucao}{10.14}\textbf{Gabarito: A.} 802.1X controla acesso baseado em porta; EAP transporta métodos de autenticação e RADIUS é backend comum.\end{solucao}
\begin{solucao}{10.15}\textbf{Gabarito: CERTO.} ICMP participa de diagnóstico e mecanismos de controle. Bloqueá-lo indiscriminadamente pode prejudicar PMTUD e troubleshooting.\end{solucao}
\begin{solucao}{10.16}\textbf{Gabarito: B.} O processo não pode prosseguir até o I/O terminar, portanto fica bloqueado/esperando.\end{solucao}
\begin{solucao}{10.17}\textbf{Gabarito: C.} Há espera circular: cada processo detém um recurso e espera pelo outro. É uma das condições clássicas de deadlock.\end{solucao}
\begin{solucao}{10.18}\textbf{Gabarito: CERTO.} \texttt{777} pode mascarar a causa do problema e concede permissões amplas. O controle correto é ajustar ownership, grupo, ACL ou privilégio de forma mínima.\end{solucao}
\begin{solucao}{10.19}\textbf{Gabarito: A.} AD depende fortemente de DNS para localizar serviços e controladores por registros apropriados.\end{solucao}
\begin{solucao}{10.20}\textbf{Gabarito: C.} RAID 5 fornece capacidade útil aproximada de $(N-1)\times S = 3\times6=18$ TB.\end{solucao}
\begin{solucao}{10.21}\textbf{Gabarito: CERTO.} Se falham dois discos de espelhos diferentes, pode haver sobrevivência; se falham os dois membros do mesmo espelho, a cópia daquele conjunto é perdida.\end{solucao}
\begin{solucao}{10.22}\textbf{Gabarito: B.} SAN oferece bloco/LUN; NAS oferece arquivo por protocolos como SMB/NFS.\end{solucao}
\begin{solucao}{10.23}\textbf{Gabarito: B.} Backup incremental guarda alterações desde o último backup de qualquer tipo. A restauração exige o último full e toda a cadeia de incrementais posteriores até o ponto desejado.\end{solucao}
\begin{solucao}{10.24}\textbf{Gabarito: CERTO.} Sem captura mais frequente de alterações, a perda potencial pode se aproximar de 24 horas, incompatível com RPO de 15 minutos.\end{solucao}
\begin{solucao}{10.25}\textbf{Gabarito: B.} Estado local prende a sessão à instância. Externalização/replicação ou desenho stateless torna escalabilidade e failover mais robustos.\end{solucao}

\begin{solucao}{10.26}\textbf{Gabarito: C.} I e II estão corretas. III está errada porque a distribuição em LACP normalmente usa hash por fluxo; um fluxo individual pode permanecer em um único membro.\end{solucao}
\begin{solucao}{10.27}\textbf{Gabarito: CERTO.} A segmentação só reduz movimento lateral quando acompanhada de políticas que limitem tráfego entre zonas.\end{solucao}
\begin{solucao}{10.28}\textbf{Gabarito: B.} Throughput TCP depende de RTT, janela, perdas e congestionamento. Banda nominal sozinha não determina taxa de um fluxo.\end{solucao}
\begin{solucao}{10.29}\textbf{Gabarito: CERTO.} MPLS VPN fornece separação lógica no backbone do provedor, não criptografia fim a fim por definição.\end{solucao}
\begin{solucao}{10.30}\textbf{Gabarito: B.} SNMPv3 melhora autenticação/privacidade, mas compartilhar credenciais e aceitar acesso de qualquer VLAN viola mínimo privilégio e segregação.\end{solucao}
\begin{solucao}{10.31}\textbf{Gabarito: B.} Os sintomas indicam pressão de memória e paginação intensa. Aumentar CPU pode pouco ajudar porque o processo está bloqueado em I/O de memória secundária.\end{solucao}
\begin{solucao}{10.32}\textbf{Gabarito: CERTO.} VM e container têm mecanismos de isolamento distintos. A VM inclui SO convidado; container compartilha kernel do host.\end{solucao}
\begin{solucao}{10.33}\textbf{Gabarito: B.} O storage compartilhado sem redundância permanece SPOF mesmo com hosts distintos.\end{solucao}
\begin{solucao}{10.34}\textbf{Gabarito: CERTO.} HA deve ser avaliada ponta a ponta, incluindo dependências compartilhadas.\end{solucao}
\begin{solucao}{10.35}\textbf{Gabarito: B.} O ponto recuperado estava dentro do RPO de 30 min, mas o serviço levou 3 h, excedendo o RTO de 1 h.\end{solucao}
\begin{solucao}{10.36}\textbf{Gabarito: C.} RAID e snapshots locais podem ser atingidos pelo mesmo comprometimento. O backup externo semanal pode deixar grande janela de perda. É preciso avaliar imutabilidade, isolamento, credenciais e testes.\end{solucao}
\begin{solucao}{10.37}\textbf{Gabarito: CERTO.} Deduplicação economiza espaço, mas não cria independência de falhas. Metadados e catálogo tornam-se componentes críticos.\end{solucao}
\begin{solucao}{10.38}\textbf{Gabarito: A.} Usar conta privilegiada em navegação e e-mail amplia exposição a phishing, malware e roubo de sessão. Contas administrativas separadas reduzem esse risco.\end{solucao}
\begin{solucao}{10.39}\textbf{Gabarito: CERTO.} Liveness deve refletir saúde do processo. Se depender de serviço externo, uma falha remota pode provocar restart storm.\end{solucao}
\begin{solucao}{10.40}\textbf{Gabarito: B.} Média global mascara períodos críticos. Percentis, picos, filas e correlação com SLA são necessários para capacidade.\end{solucao}

\begin{solucao}{10.41}\textbf{Gabarito: ERRADO.} Se o storage for perdido, snapshots locais também podem ser perdidos. Backup externo mensal tampouco sustenta RPO de 15 min. O cenário não demonstra o objetivo declarado.\end{solucao}
\begin{solucao}{10.42}\textbf{Gabarito: B.} As VLANs existem, mas a política permissiva neutraliza o controle de contenção. O achado deve focar segmentação efetiva e movimento lateral.\end{solucao}
\begin{solucao}{10.43}\textbf{Gabarito: A.} Contas nominativas, PAM/MFA, separação de privilégio e logs fora do alcance do administrador aumentam accountability e reduzem abuso.\end{solucao}
\begin{solucao}{10.44}\textbf{Gabarito: CERTO.} Diversidade aparente pode esconder mesma causa comum. Provedor, duto e ponto de entrada precisam ser avaliados.\end{solucao}
\begin{solucao}{10.45}\textbf{Gabarito: A.} Imagem sem assinatura prejudica origem/integridade; \texttt{latest} prejudica reprodutibilidade; privilégio padrão aumenta blast radius.\end{solucao}
\begin{solucao}{10.46}\textbf{Gabarito: CERTO.} Quem é auditado consegue destruir a própria evidência. Logs centralizados/protegidos e segregação são controles essenciais.\end{solucao}
\begin{solucao}{10.47}\textbf{Gabarito: B.} As duas VMs compartilham host e storage local, portanto falha física comum derruba compute e dados.\end{solucao}
\begin{solucao}{10.48}\textbf{Gabarito: B.} Job concluído indica que a tarefa terminou, não que a restauração funciona. Teste real com integridade e tempo é evidência mais forte.\end{solucao}
\begin{solucao}{10.49}\textbf{Gabarito: CERTO.} Hash demonstra integridade da cópia a partir do momento de cálculo, não veracidade ou autenticidade do conteúdo original.\end{solucao}
\begin{solucao}{10.50}\textbf{Gabarito: B.} O SLA de média não captura a cauda. O auditor deve correlacionar p95/p99 com jornadas críticas, filas e saturação antes de prescrever capacidade.\end{solucao}

